\chapter{Names Without Names: Combinators and the Matrix Machine}
\label{ch:combinators}

\section{Status, first}

The previous chapter reported on a store that has run in production for years. This one
reports on a compiler whose first build passed its test suite in the autumn of 2026, on a
software graphics driver, and whose performance has not been measured on a single program
anyone wanted to run. The reader is asked to hold that asymmetry in mind throughout, and
the chapter will keep saying so, because the argument it makes is of a kind that invites
overstatement: it is an argument that a theoretical property of the calculus ---
reflection, the thing that makes rho rho --- is exactly what makes the calculus cheap on
hardware the store was not designed for. The argument is sound as far as it goes. How far
it goes is a number nobody has yet produced.

Here is the shape of what follows. A name-free combinator calculus for rho exists, and the
version of it that was published is broken in an instructive way; repairing it requires
adding something the published version lacked, and the repaired calculus compiles rho in
linear size where the classical method is exponential. A combinator state is a multiset of
fixed-width tuples, so it is a vector; the rewrite rules are a matrix; and the one rule
that gives the calculus its computational power is also the one rule that cannot be
written down in advance and the one that breaks the matrix's triangular structure. The
feature that makes the theory hard makes the machine cheap: reflection is a lookup table.
And the speedup available from running this on an accelerator is the \emph{width} of the
computation and nothing else --- a falsifiable thesis that should be tested before anyone
builds more hardware than has already been built.

\section{A calculus with no bound names}
\label{cb:sec:calculus}

Rho has no name-restriction operator: names are quoted processes, and freshness comes
from quotation rather than from a binder. It nonetheless has one binder, the input
prefix. Yoshida showed how to eliminate the input binder of the asynchronous
$\pi$-calculus in favor of combinators, at the cost of keeping $\mathsf{new}$. The
combination is suggestive: a calculus with no $\mathsf{new}$, compiled by a discipline
that removes the remaining binder, would be a concurrency calculus with no bound names
whatsoever. That is what the rho combinators proposed~\cite{rhocomb}, and the current
draft of the work~\cite{rhocomb3} asks whether the proposal works.

The atoms are nine. Six are \emph{routers}, which move a message without reading it:
$\mathsf{d}(a,b,c)$ duplicates a message at $a$ to $b$ and $c$; $\mathsf{k}(a)$ kills
one; $\mathsf{fw}(a,b)$ forwards $a$ to $b$; $\mathsf{bl}(a,b)$ and $\mathsf{br}(a,b)$
bind the received payload as the left or right argument of a new forwarder; and
$\mathsf{sy}(a,b,c)$ synchronizes, emitting a forwarder from $b$ to $c$ when a message
arrives at $a$. One is the \emph{destructor} $\mathsf{ev}(a)$, which consumes a message at
$a$ and \emph{runs its payload}. Two carry values: the message $\mathsf{m}(a,v)$ and the
store $\mathsf{q}(a,p)$, a message that only a forwarder can read. The rules are two-atom
rewrites, each a message meeting one router or the destructor at a shared name.

\begin{figure}[t]
\hrule\vspace{4pt}
\[
\begin{array}{l@{\qquad}l}
 \mathsf{d}(a,b,c) \mid \mathsf{m}(a,v) \rewrite \mathsf{m}(b,v) \mid \mathsf{m}(c,v)
   & \mathsf{bl}(a,b) \mid \mathsf{m}(a,v) \rewrite \mathsf{fw}(v,b)\\
 \mathsf{k}(a) \mid \mathsf{m}(a,v) \rewrite \nil
   & \mathsf{br}(a,b) \mid \mathsf{m}(a,v) \rewrite \mathsf{fw}(b,v)\\
 \mathsf{fw}(a,b) \mid \mathsf{m}(a,v) \rewrite \mathsf{m}(b,v)
   & \mathsf{sy}(a,b,c) \mid \mathsf{m}(a,v) \rewrite \mathsf{fw}(b,c)\\
 \mathsf{ev}(a) \mid \mathsf{m}(a,v) \rewrite *v
   & \mathsf{fw}(a,b) \mid \mathsf{q}(a,p) \rewrite \mathsf{m}(b,@p)\\[4pt]
 \multicolumn{2}{l}{\mathsf{cons}_{\mid}(a,b,c) \mid \mathsf{m}(a,@p) \mid \mathsf{m}(b,@r)
   \rewrite \mathsf{m}(c,@(p \mid r))}\\
 \multicolumn{2}{l}{\mathsf{cons}_{A}(a_{1},\dots,a_{n},f) \mid \mathsf{m}(a_{1},v_{1})
   \mid \cdots \mid \mathsf{m}(a_{n},v_{n})
   \rewrite \mathsf{m}(f,@A(v_{1},\dots,v_{n}))}
\end{array}
\]
\vspace{2pt}\hrule
\caption{The rules of the repaired combinator calculus, closed under $\mid$ and the
equations of a commutative monoid on $\mid$. The last line is a schema, one rule per
atom shape $A$ a program needs.}
\label{cb:fig:rules}
\end{figure}

\paragraph{One equation breaks the published calculus.} The published combinators inherit
$*@P \equiv P$ from a presentation in which quote and drop are mutually inverse, and they
write the destructor and the drop with the same symbol. In the combinator setting the
destructor is not a process former but a guard; and since every name is $@Q$ for some $Q$,
identifying the guard with the drop makes an arbitrary running process match the guard.

\begin{proposition}[The defect]
\label{cb:prop:defect}
If the drop and the destructor are one constructor and $*@p = p$, then for all $P,Q$:
$\;Q \mid \mathsf{m}(@Q,@P) \rewrite P$.
\end{proposition}

\begin{proof}
Every name is $@Q$ for some $Q$, so under the conflation $\mathsf{ev}(@Q) = Q$. Reduction
is closed under the equations, so the redex $\mathsf{ev}(a) \mid \mathsf{m}(a,@P)$ with
$a=@Q$ matches $Q \mid \mathsf{m}(@Q,@P)$ and fires.
\end{proof}

An arbitrary running process can be annihilated and replaced by an arbitrary payload by
anyone who knows its quotation. No faithfulness theorem survives this. The repair is to
separate the drop from the guard --- Figure~\ref{cb:fig:rules} does --- which leaves the
equation optional; the source gives both presentations.

\paragraph{Reflection was half-implemented, and constructors complete it.} The rule
$\mathsf{ev}(a) \mid \mathsf{m}(a,@P) \rewrite P$ opens a name that arrives at run time.
In the published calculus nothing \emph{builds} one: combinator reduction never produces
a name not already present in the redex, hereditarily. Rho reduction does, at every step,
because substituting under a quotation turns $@(*y \mid R)$ into $@(Q \mid R)$ --- a name in
neither participant. So the published calculus cannot encode rho at all. The missing
operation is the construction of a name from names, and it is supplied by
\emph{constructors}: one for parallel composition, $\mathsf{cons}_{\mid}$, and one per atom
shape, each taking its arguments as messages and emitting the quotation of the term built
from them. With constructors there is no name-opaque fragment to restrict to. They are the
combinator-level image of the one thing the Pledge said reflection had to do cheaply: let
a network of agents be the sort of thing that can be named.

\paragraph{The storage combinator is about opacity, not power.} $\mathsf{q}(a,p)$, released
by a forwarder rather than by a message, adds nothing to expressiveness. What it adds is
\emph{silence}: it is a message that only a redirection can read, which is what stored code
and intermediate results of name construction need if the encoding is to be opaque to
observers.

\paragraph{Prefix elimination is linear.} Reflection allows a continuation to be quoted
once and released by a three-atom gate. Yoshida must instead push the prefix through the
body, splitting at every parallel composition and delaying every atom, which costs a
constant factor per level of nesting.

\begin{proposition}[Succinctness]
\label{cb:prop:blowup}
On an explicit family of terms of nesting depth $d$, Yoshida's elimination produces
$(3^{d-1}+1)/2$ atoms where the reflective encoding produces $3d+1$. In general the image
of a rho term $P$ under the two-phase translation has $O(|P|)$ atoms.
\end{proposition}

That is the first payoff, and it is a theorem~\cite[Prop.~8.1]{rhocomb3}: \emph{storing a
continuation as a name makes compiling away the input binder linear in program size, where
the classical method blows up exponentially in nesting depth.}

\paragraph{Correctness is full abstraction relative to encoded contexts.} Observation is by
context-labeled bisimilarity, with the observer class taken to be the image of the
translation on source contexts, and the theorem is that two rho terms are bisimilar iff
their images are bisimilar against encoded observers~\cite[\S7]{rhocomb3}. The
restriction to encoded observers is forced rather than convenient --- an unrestricted
combinator context can see the apparatus --- by the same argument that forces it for the
encoding of the $\pi$-calculus into rho (Chapter~\ref{ch:obsext}). The proofs are at the
level of sketches; the source says so.

\paragraph{An expressiveness lattice.} Including or excluding combinators gives a lattice
of fragments, and three invariants draw it. Without the destructor the calculus is
strongly normalizing and so not Turing complete; without constructors it cannot build a
name; and the duplicator $\mathsf{d}$ is conjectured necessary for divergence. The
constructor axis is Boolean. Whether the three binders $\mathsf{bl},\mathsf{br},\mathsf{sy}$
are separable from one another is open. The lattice has two orderings, expressive and
observational, and they do not coincide, because the calculus has no hiding. We mention it
because the next section will show that the same three invariants are the three finiteness
conditions of a matrix.

\section{States are vectors, rules are a matrix}
\label{cb:sec:matrix}

Everything in this section is over the Booleans or the natural numbers; the weighted and
complex readings of Chapter~\ref{ch:grading} attach to the same operator later and nothing
here depends on them. The source is a design analysis~\cite{rhocombmat}, and it marks
each of its claims as settled, as a route believed to work, as open, or as a gate on the
work downstream; we carry the markings.

\paragraph{What a graphics processor rewards.} An accelerator is a bulk-synchronous machine
with a very wide, very shallow execution model. It rewards flat, indexable state --- arrays
of fixed-width records at predictable addresses, pointer chasing being the single most
expensive thing one can ask of it; uniform work across lanes --- threads that take the same
branch; bulk-synchronous phases --- one kernel does one thing to everything, then a
barrier; and standard sparse primitives --- sparse matrix products, segmented scans, sorts
and hash tables, which inherit decades of tuning. An RSpace-style interpreter does
approximately the opposite of all four. Its state is a store of variably shaped terms
reached by pointer; matching is a recursive walk over nested syntax; conflicts are
resolved by locks; and the inner loop is a hash lookup followed by unification. None of
that is a criticism of RSpace, which is a good design for a general machine. It is an
itemized statement of the gap a compiler standing between the two has to close.

Four facts about the repaired combinators do the closing. \emph{Atoms are flat}: every atom
is a head symbol applied to at most $w$ arguments, $w$ known at compile time, and every
argument is a name --- there is no nesting, no binder, no pattern variable. \emph{The
multiset is complete}: the multiset of top-level atoms of a drop-free normal form is a
complete invariant for structural congruence. \emph{Rules are small and subject-joined}:
two to $w$ premises, each a single atom, sharing names. \emph{Name equality is decidable and
internable}: on closed normal forms it is decided by normalization and multiset comparison,
and quotation is injective there, so names can be hash-consed to integers once and
compared by integer equality thereafter.

\begin{definition}[Species, and the state space]
\label{cb:def:species}
A \emph{species} is a head symbol together with a tuple of names of the matching arity. A
\emph{state} is a finite multiset over species --- an \emph{occupation vector}
$\gamma : \mathrm{Spec} \to \mathbb{N}$ of finite support. The state space is the free
module on such $\gamma$; equivalently, the free commutative algebra generated by the
species, a state being a monomial $\prod_\sigma (a^{\dagger}_\sigma)^{\gamma(\sigma)}
\lvert\nil\rangle$.
\end{definition}

\begin{principle}[Structural congruence is free]
\label{cb:prin:free}
The equations of the calculus are not imposed on the representation and are not discharged
by a normalization pass. Commutativity, associativity and the unit of $\mid$ are the
algebra of the multiset; $*@p = p$ and $@*a = a$ are discharged once, at interning time.
No run-time cost.
\end{principle}

This is the first place the combinators pay. In a direct rho machine, matching modulo
structural congruence is a recurring cost and a recurring source of subtlety; here it is
a property of the data structure. It is also, the reader will notice, Chapter~\ref{ch:rspace}'s
Proposition~\ref{rs:prop:hom} said a second way: the store made $\Par$ into multiset union
of bags at keys; the matrix machine makes the whole state into one multiset over species,
and the keys have become coordinates.

\begin{definition}[The rule operator]
\label{cb:def:operator}
For a species $\sigma$ let $a^{\dagger}_\sigma$ add one occurrence and $a_\sigma$ remove
one, with $a_\sigma \lvert\gamma\rangle = 0$ when $\gamma(\sigma) = 0$. Each rule of
Figure~\ref{cb:fig:rules} is a monomial in these operators --- annihilate the premises,
create the conclusions --- schematic in the names; the \emph{transfer operator} is the sum
over rules and name instances, $T = \sum_r H_r$.
\end{definition}

\begin{proposition}[Matrix elements count derivations]
\label{cb:prop:massaction}
$\langle\gamma'\lvert T \rvert\gamma\rangle$ is the number of distinct rule instances
carrying $\gamma$ to $\gamma'$, which for a rule with premises of multiplicities
$\gamma(\sigma_1),\dots,\gamma(\sigma_k)$ is the product of those multiplicities.
\end{proposition}

The entries are mass-action counts. A combinator program is, without any choice having
been made, a chemical reaction network: species, reactions, and a transfer operator whose
entries are the number of ways each reaction can fire. Chapter~\ref{ch:grading}'s
multiplicity factor --- count redexes by occurrence, or run the system at half rate --- is
here the definition of the matrix element. The weighted reading attaches at exactly one
place, a diagonal factor in each rule's operator term, which is why the Gillespie reading
of that chapter and the matrix machine of this one are one object.

\paragraph{The self-indexing problem.} Species are indexed by names, and names are quoted
processes, so by completeness of the multiset they are finite multisets over species. The
index set of the basis therefore satisfies
\[
  \mathrm{Nm} \;\cong\; M_{\mathrm{fin}}(\mathrm{Spec}), \qquad
  \mathrm{Spec} \;=\; \textstyle\coprod_{A} \mathrm{Nm}^{\mathrm{ar}(A)} .
\]
The basis of the state space indexes itself. This is reflection seen through the linear
algebra, and it is the reason no finite matrix can be written down in advance for the full
calculus. The solution of the domain equation is the initial algebra: names are finite
trees, so the name set is countable, finitely branching, and graded by depth of quotation.

\begin{definition}[The intern table]
\label{cb:def:intern}
The machine realizes the name set as a bijection between integers and canonical component
multisets, in two directions: $\mathsf{quote}$, a hash-cons --- an \emph{encode}; and
$\mathsf{drop}$, a table read --- a \emph{decode}.
\end{definition}

\begin{theorem}[The lattice classifies the matrix]
\label{cb:thm:three}
Let $T$ be a combinator term over a constructor set $S$.
\begin{enumerate}[label=(\roman*),leftmargin=2.4em]
\item \emph{(Index set; constructor axis.)} Every name occurring in a reduct of $T$ lies in
the closure of $T$'s names under the node shapes in $S$. If $S=\emptyset$ the reachable
species set is finite; if not, it is the free $S$-algebra on $T$'s names, countable and
graded by depth.
\item \emph{(Nilpotency; destructor axis.)} Every rule but $\mathsf{ev}$ strictly decreases
a weight counting non-message atoms. Without $\mathsf{ev}$ the reachable operator is
strictly block upper triangular with respect to that grading and $T^{2N+1} = 0$ for $N$ the
number of non-message atoms.
\item \emph{(Occupation bound; duplication axis.)} Without $\mathsf{d}$ and without
$\mathsf{ev}$, the combined occupation of message and store species is non-increasing.
\end{enumerate}
\end{theorem}

The three invariants the expressiveness lattice was drawn with are, without modification,
the three finiteness conditions a matrix representation needs, one per axis of the lattice.
That is the source's main structural observation and it is proved. The corollary the
original proposal wanted --- that the reduction graph is an upper triangular matrix, so
that reachability is a finite Neumann series computed by repeated squaring --- holds
\emph{exactly on the destructor-free fragment}. With $\mathsf{ev}$ the graph can have
cycles, because reflection admits terms that step to a structurally congruent copy of
themselves, and by Principle~\ref{cb:prin:free} congruent means equal.

\begin{remark}[One atom, three characterizations]
\label{cb:rmk:evthrice}
The destructor $\mathsf{ev}$ is, simultaneously: the carrier of Turing power, by the
strong-normalization result for the fragment without it; the unique breaker of the
grading, by Theorem~\ref{cb:thm:three}(ii); and the unique rule whose operator cannot be
written schematically in advance, because its conclusion is whatever the decoded name
says. Three independent characterizations of one atom, from logic, from order theory and
from implementation. The remark in the combinator paper that ``the destructor carries the
computational power that reflection supplies'' acquires an operational reading: the power
is exactly the ability to turn an index into a state.
\end{remark}

The recommendation that follows is the one the source makes, and it reverses the original
proposal. For running a program one needs a \emph{frontier}, not a closure: the current
state and its successors. Triangularity is a property of the whole graph and buys nothing
at the frontier; a general sparse non-negative operator applied to a sparse vector is
entirely adequate, and iterative methods are indifferent to cycles. The upper triangular
matrix is the right object for model checking --- or, working in the unfolding, for the
history monad's universal cover, which is acyclic by construction --- and the wrong object
for execution. \emph{Do not ask for it.}

\section{What actually gets multiplied}
\label{cb:sec:multiply}

Not the reduction graph. The object the machine multiplies is the sparse \emph{incidence
matrix} between atom occurrences and names --- term-sized, not graph-sized --- and three
operations on it are the whole of a step. \emph{Redex finding is a join}: the rules are
subject-joined, so the instances of a rule enabled in the current state are found by
joining the atom table with itself on the shared name column, which is a sparse product.
\emph{Conflict detection is a second product}: two redexes conflict when they consume a
common atom occurrence, and the conflict graph is $BB^{\top}$ for $B$ the redex-atom
incidence. \emph{Independent-set selection is a standard parallel primitive}. Each step
fires a \emph{maximal independent set} of redexes, which is one barrier per causal layer
--- and firing a maximal independent set per step is sound with respect to the
interleaving semantics, so that machine steps count \emph{causal depth}. This is the
contention-component structure of Chapter~\ref{ch:grading} with the components found by
linear algebra: the independent-set step is the $k$-ary generalization of the
contention-set normalization, and the source files reconciling the two as an obligation.

\begin{principle}[Reflection is an intern table]
\label{cb:prin:reflection}
The decode rule reads the intern table at $v$ and appends the stored component multiset to
the state: a table lookup followed by a bulk copy of a contiguous block, among the cheapest
things an accelerator does. The constructors build a small canonical term and hash-cons it:
a concurrent insert-or-get, a standard device primitive. The feature that makes the
calculus theoretically hard --- that a name is a program and a program can be recovered
from a name --- is, at the level of the machine, exactly a bidirectional intern table. The
theoretical difficulty and the implementation difficulty do not coincide here, and they
point in opposite directions: reflection is what obstructs a static matrix and what makes
the dynamic one cheap, because storing a continuation as an index is what keeps atoms
fixed-width.
\end{principle}

This also explains Proposition~\ref{cb:prop:blowup} from the machine side. Yoshida's
elimination must rewrite a continuation to delay it; the reflective encoding stores it as a
name and pays for a gate. A machine that stores continuations as integers in a table is
doing the same thing, and the $O(|P|)$ bound is what keeps the state array small enough
to be worth the transfer to a device. The intern table is the one place the
bulk-synchronous discipline is violated --- it is shared mutable state --- but both of its
operations are commutative and idempotent in effect, so it can be updated by an unordered
batch insert per step rather than by locking. Interning order is a gauge freedom: two runs
may assign different integers to the same name, the states are equal up to the canonical
bijection, and the machine must never let an integer comparison stand in for a semantic
distinction between runs.

\section{What has been measured, and the thesis}
\label{cb:sec:measured}

A reference implementation accompanies the design note, and its measurements are step
counts and widths, not timings. On four synthetic term families with analytically known
widths --- $k$ independent forwarder chains of length $d$; $c$ messages and $c$ consumers
on one channel; a binary duplicator tree of height $h$; a construct-then-eval family ---
the number of parallel steps equals the causal depth exactly and the ratio of sequential
to parallel steps equals the width exactly: $64.00$ for $64$ chains, $16.00$ for $16$
contenders, $18.14$ for a duplicator tree of height six. The join formulation of redex
finding agrees with naive enumeration on $400$ of $400$ random states.

\begin{principle}[The width thesis]
\label{cb:prin:width}
The speedup available from this route is the mean width of the reduction, and nothing
else. Sequential terms gain nothing; the machine adds a constant per-step overhead of two
sparse products and an independent-set selection. The thesis is falsifiable by measuring
the width distribution of compiled rho programs, and should be falsified or confirmed
before further device work is undertaken.
\end{principle}

Three things are not claimed, and the source says so in its own voice: that any of this
beats a well-tuned sequential interpreter on real workloads; that the reduction graph should
ever be materialized; that maximal progress is the right resolver rather than \emph{a}
resolver whose cost is an open question. And one thing is marked as the gate on the whole
enterprise: \emph{the width of the combinator image of a compiled rho program has not been
measured.} The measurements establish that the construction is correct and that step
count is causal depth on families whose depth is known. They establish nothing about real
programs, because that measurement requires the compiler, not the machine.

\section{The compiler, as it stands}
\label{cb:sec:saphire}

The compiler is called Saph1r3. It is a Rust workspace of eight crates, about eight thousand
lines, whose front end is the kernel-rho parser and normalizer of the campf1r3 project
taken unmodified, and whose pipeline is: parse and normalize; \emph{phase one}, a
homomorphism on the normal form into the combinators extended with a binder; \emph{phase
two}, which eliminates the binder by giving each run-time instance an address and erecting
its apparatus under that address with an instantiation combinator; then deployment to a
host matrix machine or to a device bundle --- a program matrix, kernels generated from the
rule table in the WebGPU shading language, and a manifest --- targeting Vulkan, Metal or
DirectX through a portable layer. The device finds redexes and selects the independent set;
the host commits each step; the intern table stays on the host.

Its test suite reports forty-eight passing tests at the time of writing: a hundred-seed
regression against a reference reducer on the paper's own adversarial term and its
variants; a nineteen-program corpus covering every occurrence kind, nesting, relays,
races and received code with inputs; device traces byte-identical to host traces; address
independence; canonicity of the encoding; golden files for the full reduction traces of
two programs; an adversarial cross-instance scheduler with zero disagreements over a hundred
seeds; and a nesting family on which phase-one atoms grow linearly in depth, which is
Proposition~\ref{cb:prop:blowup} observed. The device tests were run on a software Vulkan
driver. No performance target has been set, let alone met.

The build also found things wrong in the specifications it implements, and reported them.
Collapsing the hole indices of nested contexts to a single marker is unsound for nested
units --- a parent's instantiation would fill a child's hole --- and de~Bruijn indices over
context nodes repair it. The specification's address supply for received code was
unstated, and the build closes the gap with one address per release site at a shared
static channel. A quoted process containing an input is a unit and must be erected like an
input body. These are the ordinary findings of a first implementation, and they are
recorded in the repository's own words as findings for the papers.

\begin{remark}[What is nascent, stated plainly]
\label{cb:rmk:nascent}
The store-backed machine --- the one that would let the matrix machine run against the
production store of Chapter~\ref{ch:rspace} rather than its own interned state --- is not
built. There is no browser build. The alternative erection scheme the logic paper favors
reaches twelve of twenty-one corpus programs and not the regression term itself. The
specification says the machine has no external dependencies and the code depends on a
graphics library. And the one number the whole route turns on --- the width distribution
of compiled rho programs --- has not been computed, although the compiler that would
compute it now exists. In the vocabulary of this book: a hypothesis has been installed as a
guard, the assay has been designed, and no message has yet arrived on the probe channel.
\end{remark}

\section{Two realizations of one rule}
\label{cb:sec:two}

The store and the matrix machine implement the same COMM rule and agree on what it means,
and they disagree about everything else. The store keys the state by channel and finds a
redex by looking up a key; the machine flattens the state to species and finds every
redex at once by a join. The store resolves a race at a key, under a lock, with a seeded
shuffle; the machine resolves every race in a layer at once, by independent-set selection,
under maximal progress. The store is a general machine's realization and runs anything;
the machine is an accelerator's realization and runs what is wide. The store carries
reflection as a content-addressed key --- a quoted process \emph{is} its hash; the machine
carries it as an intern table --- a quoted process \emph{is} its integer. And the store
knows the two messages at $a$ are two because it keeps a bag; the machine knows it because
the matrix element is $2$.

Which is to say that the resource layer of Chapter~\ref{ch:rspace} --- the layer the
behavioral layer forgets, the layer that keeps counts and keys --- is the layer both
machines live in, and that they are two coordinate systems on it. A reader who has
followed the book's habit of asking what a construction prices will want to know what the
change of coordinates costs and buys. That is the next chapter.

% ===================================================================
